\documentclass[11pt,a4paper]{article}
\usepackage[utf8]{inputenc}
\usepackage[spanish]{babel}
\usepackage{amsmath,amssymb,amsthm}
\usepackage{graphicx}
\usepackage{xcolor}
\usepackage{geometry}
\usepackage{fancyhdr}
\usepackage{listings}
\usepackage{hyperref}
\usepackage{booktabs}
\usepackage{float}

\geometry{margin=2.5cm}

\definecolor{codegreen}{rgb}{0,0.6,0}
\definecolor{codegray}{rgb}{0.5,0.5,0.5}
\definecolor{codepurple}{rgb}{0.58,0,0.82}
\definecolor{backcolour}{rgb}{0.95,0.95,0.92}

\lstdefinestyle{pythonstyle}{
    backgroundcolor=\color{backcolour},
    commentstyle=\color{codegreen},
    keywordstyle=\color{blue},
    numberstyle=\tiny\color{codegray},
    stringstyle=\color{codepurple},
    basicstyle=\ttfamily\footnotesize,
    breakatwhitespace=false,
    breaklines=true,
    captionpos=b,
    keepspaces=true,
    numbers=left,
    numbersep=5pt,
    showspaces=false,
    showstringspaces=false,
    showtabs=false,
    tabsize=2
}

\lstset{style=pythonstyle}

\pagestyle{fancy}
\fancyhf{}
\fancyhead[L]{\small Métodos de Análisis de Datos 1}
\fancyhead[R]{\small Unidad 3 - Regresión y Correlación}
\fancyfoot[C]{\thepage}

\title{\textbf{Unidad 3: Modelado Estadístico}\\ \Large Regresión y Correlación}
\author{Métodos de Análisis de Datos 1\\Universidad Nacional del Sur}
\date{Primer Semestre 2026}

\begin{document}

\maketitle
\tableofcontents
\newpage

%=============================================================================
\section{Introducción al Modelado Estadístico}
%=============================================================================

El modelado estadístico es el proceso de construir representaciones matemáticas de fenómenos del mundo real utilizando datos observados. En esta unidad nos centraremos en modelos de \textbf{regresión}, que permiten estudiar y cuantificar relaciones entre variables.

\subsection{¿Qué es un modelo de regresión?}

Un modelo de regresión describe cómo una variable respuesta $Y$ (también llamada variable dependiente) se relaciona con una o más variables predictoras $X_1, X_2, \ldots, X_p$ (variables independientes o explicativas).

\textbf{Objetivos principales:}
\begin{itemize}
    \item \textbf{Descripción:} Entender cómo se relacionan las variables
    \item \textbf{Predicción:} Estimar valores de $Y$ para nuevos valores de $X$
    \item \textbf{Inferencia:} Determinar qué variables son importantes y cuantificar su efecto
    \item \textbf{Control:} Ajustar por variables confusoras en estudios observacionales
\end{itemize}

\subsection{Tipos de modelos de regresión}

\begin{table}[H]
\centering
\begin{tabular}{ll}
\toprule
\textbf{Tipo} & \textbf{Características} \\
\midrule
Regresión lineal simple & 1 predictor, relación lineal \\
Regresión lineal múltiple & $\geq 2$ predictores, relación lineal \\
Regresión polinomial & 1 predictor, relación no lineal \\
Regresión no paramétrica & Relaciones complejas sin forma funcional \\
Regresión logística & Variable respuesta binaria \\
Regresión de Poisson & Variable respuesta de conteo \\
\bottomrule
\end{tabular}
\end{table}

En esta unidad nos concentraremos en modelos de \textbf{regresión lineal}.

%=============================================================================
\section{Regresión Lineal Simple}
%=============================================================================

\subsection{Definición del modelo}

La regresión lineal simple modela la relación entre una variable respuesta $Y$ y un predictor $X$ mediante una línea recta:

\begin{equation}
Y_i = \beta_0 + \beta_1 X_i + \epsilon_i, \quad i = 1, 2, \ldots, n
\end{equation}

donde:
\begin{itemize}
    \item $Y_i$: valor observado de la variable respuesta para la observación $i$
    \item $X_i$: valor del predictor para la observación $i$
    \item $\beta_0$: ordenada al origen (intercepto)
    \item $\beta_1$: pendiente (cambio en $Y$ por unidad de cambio en $X$)
    \item $\epsilon_i$: error aleatorio, $\epsilon_i \sim N(0, \sigma^2)$ independientes
    \item $n$: tamaño de la muestra
\end{itemize}

\subsection{Supuestos del modelo}

Para que la inferencia sea válida, se asumen:

\begin{enumerate}
    \item \textbf{Linealidad:} La relación entre $X$ e $Y$ es lineal
    \item \textbf{Independencia:} Las observaciones son independientes entre sí
    \item \textbf{Homocedasticidad:} La varianza de los errores es constante: $\text{Var}(\epsilon_i) = \sigma^2$ para todo $i$
    \item \textbf{Normalidad:} Los errores siguen una distribución normal: $\epsilon_i \sim N(0, \sigma^2)$
    \item \textbf{No multicolinealidad:} (más relevante en regresión múltiple)
\end{enumerate}

\subsection{Estimación de parámetros: Mínimos Cuadrados Ordinarios (MCO)}

El método de \textbf{mínimos cuadrados ordinarios} (MCO, o OLS en inglés) estima $\beta_0$ y $\beta_1$ minimizando la suma de los cuadrados de los residuos:

\begin{equation}
\text{SCR} = \sum_{i=1}^{n} (Y_i - \hat{Y}_i)^2 = \sum_{i=1}^{n} (Y_i - \hat{\beta}_0 - \hat{\beta}_1 X_i)^2
\end{equation}

Las estimaciones por MCO son:

\begin{align}
\hat{\beta}_1 &= \frac{\sum_{i=1}^{n} (X_i - \bar{X})(Y_i - \bar{Y})}{\sum_{i=1}^{n} (X_i - \bar{X})^2} = \frac{S_{XY}}{S_{XX}} \\
\hat{\beta}_0 &= \bar{Y} - \hat{\beta}_1 \bar{X}
\end{align}

donde $\bar{X}$ y $\bar{Y}$ son las medias muestrales de $X$ e $Y$.

\subsection{Interpretación de los coeficientes}

\begin{itemize}
    \item \textbf{Intercepto} $\hat{\beta}_0$: Valor esperado de $Y$ cuando $X = 0$
    \item \textbf{Pendiente} $\hat{\beta}_1$: Cambio esperado en $Y$ cuando $X$ aumenta en 1 unidad
\end{itemize}

\textbf{Ejemplo:} Si modelamos el precio de casas ($Y$, en miles de \$) en función de los metros cuadrados ($X$) y obtenemos:
\[
\hat{Y} = 50 + 1.2 X
\]
entonces:
\begin{itemize}
    \item $\hat{\beta}_0 = 50$: Una casa de 0 m$^2$ tendría un precio base de 50 mil \$ (interpretación sin sentido práctico en este caso)
    \item $\hat{\beta}_1 = 1.2$: Cada m$^2$ adicional aumenta el precio esperado en 1.2 mil \$ (1200 \$)
\end{itemize}

\subsection{Coeficiente de determinación $R^2$}

El $R^2$ mide la proporción de variabilidad de $Y$ explicada por el modelo:

\begin{equation}
R^2 = \frac{\text{SCE}}{\text{SCT}} = 1 - \frac{\text{SCR}}{\text{SCT}}
\end{equation}

donde:
\begin{itemize}
    \item $\text{SCT} = \sum_{i=1}^{n} (Y_i - \bar{Y})^2$: suma de cuadrados total
    \item $\text{SCE} = \sum_{i=1}^{n} (\hat{Y}_i - \bar{Y})^2$: suma de cuadrados explicada (por el modelo)
    \item $\text{SCR} = \sum_{i=1}^{n} (Y_i - \hat{Y}_i)^2$: suma de cuadrados residual
\end{itemize}

\textbf{Propiedades:}
\begin{itemize}
    \item $0 \leq R^2 \leq 1$
    \item $R^2 = 0$: el modelo no explica nada ($X$ no ayuda a predecir $Y$)
    \item $R^2 = 1$: ajuste perfecto (todos los puntos sobre la recta)
    \item En regresión lineal simple: $R^2 = r_{XY}^2$ (cuadrado de la correlación de Pearson)
\end{itemize}

\textbf{Advertencia:} Un $R^2$ alto no implica causalidad ni garantiza que el modelo sea adecuado.

\subsection{Inferencia sobre los coeficientes}

Bajo los supuestos del modelo, los estimadores $\hat{\beta}_0$ y $\hat{\beta}_1$ siguen distribuciones normales:

\begin{align}
\hat{\beta}_1 &\sim N\left(\beta_1, \frac{\sigma^2}{S_{XX}}\right) \\
\hat{\beta}_0 &\sim N\left(\beta_0, \sigma^2 \left[\frac{1}{n} + \frac{\bar{X}^2}{S_{XX}}\right]\right)
\end{align}

En la práctica, $\sigma^2$ se estima con:
\begin{equation}
s^2 = \frac{\text{SCR}}{n - 2}
\end{equation}

\textbf{Test de hipótesis:} Para probar $H_0: \beta_1 = 0$ vs $H_1: \beta_1 \neq 0$:

\begin{equation}
t = \frac{\hat{\beta}_1}{\text{SE}(\hat{\beta}_1)} \sim t_{n-2}
\end{equation}

donde $\text{SE}(\hat{\beta}_1) = \sqrt{s^2 / S_{XX}}$.

Si $|t| > t_{n-2, \alpha/2}$, rechazamos $H_0$ al nivel $\alpha$.

\textbf{Intervalos de confianza:}
\begin{equation}
\text{IC}_{1-\alpha}(\beta_1) = \hat{\beta}_1 \pm t_{n-2, \alpha/2} \cdot \text{SE}(\hat{\beta}_1)
\end{equation}

\subsection{Predicción}

Dado un nuevo valor $X_0$, podemos calcular:

\textbf{1. Predicción puntual:}
\begin{equation}
\hat{Y}_0 = \hat{\beta}_0 + \hat{\beta}_1 X_0
\end{equation}

\textbf{2. Intervalo de confianza para la media} $E(Y | X = X_0)$:
\begin{equation}
\hat{Y}_0 \pm t_{n-2, \alpha/2} \cdot s \sqrt{\frac{1}{n} + \frac{(X_0 - \bar{X})^2}{S_{XX}}}
\end{equation}

\textbf{3. Intervalo de predicción para una nueva observación:}
\begin{equation}
\hat{Y}_0 \pm t_{n-2, \alpha/2} \cdot s \sqrt{1 + \frac{1}{n} + \frac{(X_0 - \bar{X})^2}{S_{XX}}}
\end{equation}

El intervalo de predicción es siempre más ancho que el de confianza porque incluye tanto la incertidumbre del modelo como la variabilidad individual.

%=============================================================================
\section{Diagnóstico del Modelo}
%=============================================================================

\subsection{Análisis de residuos}

Los \textbf{residuos} son las diferencias entre valores observados y predichos:
\begin{equation}
e_i = Y_i - \hat{Y}_i
\end{equation}

Examinamos los residuos para verificar los supuestos del modelo.

\subsubsection{Gráficos de diagnóstico principales}

\textbf{1. Residuos vs Valores Ajustados:}
\begin{itemize}
    \item \textbf{Qué buscar:} Patrón aleatorio sin estructura
    \item \textbf{Problemas detectables:} No linealidad, heterocedasticidad
    \item Si hay patrón en forma de embudo $\Rightarrow$ heterocedasticidad
    \item Si hay curvatura $\Rightarrow$ relación no lineal
\end{itemize}

\textbf{2. Q-Q Plot (Normal):}
\begin{itemize}
    \item \textbf{Qué buscar:} Puntos sobre la diagonal
    \item \textbf{Problemas detectables:} Violación de normalidad
\end{itemize}

\textbf{3. Scale-Location (Raíz de residuos estandarizados vs valores ajustados):}
\begin{itemize}
    \item \textbf{Qué buscar:} Banda horizontal con dispersión uniforme
    \item \textbf{Problemas detectables:} Heterocedasticidad
\end{itemize}

\textbf{4. Residuos vs Leverage (Distancia de Cook):}
\begin{itemize}
    \item \textbf{Qué buscar:} Identificar observaciones influyentes
    \item \textbf{Problemas detectables:} Outliers con alta influencia
\end{itemize}

\subsection{Tests formales}

\textbf{Test de normalidad:}
\begin{itemize}
    \item Shapiro-Wilk: $H_0$: los residuos siguen una distribución normal
\end{itemize}

\textbf{Test de homocedasticidad:}
\begin{itemize}
    \item Breusch-Pagan: $H_0$: varianza constante
    \item White: versión más robusta
\end{itemize}

\subsection{Medidas de influencia}

\textbf{Leverage (apalancamiento):} Mide qué tan lejos está $X_i$ del promedio $\bar{X}$
\begin{equation}
h_i = \frac{1}{n} + \frac{(X_i - \bar{X})^2}{S_{XX}}
\end{equation}

\textbf{Distancia de Cook:} Combina residuo y leverage para medir influencia
\begin{equation}
D_i = \frac{(e_i)^2}{2s^2} \cdot \frac{h_i}{1 - h_i}
\end{equation}

Observaciones con $D_i > 1$ son influyentes y deben ser examinadas cuidadosamente.

%=============================================================================
\section{Regresión Lineal Múltiple}
%=============================================================================

\subsection{Modelo con múltiples predictores}

La regresión lineal múltiple extiende el modelo simple a $p$ predictores:

\begin{equation}
Y_i = \beta_0 + \beta_1 X_{i1} + \beta_2 X_{i2} + \cdots + \beta_p X_{ip} + \epsilon_i
\end{equation}

En forma matricial:
\begin{equation}
\mathbf{Y} = \mathbf{X} \boldsymbol{\beta} + \boldsymbol{\epsilon}
\end{equation}

donde:
\begin{itemize}
    \item $\mathbf{Y}$: vector de respuestas ($n \times 1$)
    \item $\mathbf{X}$: matriz de diseño ($n \times (p+1)$), incluye columna de 1's para el intercepto
    \item $\boldsymbol{\beta}$: vector de coeficientes ($(p+1) \times 1$)
    \item $\boldsymbol{\epsilon}$: vector de errores ($n \times 1$)
\end{itemize}

\subsection{Estimación por MCO}

El estimador de MCO en forma matricial:
\begin{equation}
\hat{\boldsymbol{\beta}} = (\mathbf{X}^\top \mathbf{X})^{-1} \mathbf{X}^\top \mathbf{Y}
\end{equation}

\textbf{Supuestos:} Los mismos que en regresión simple, más:
\begin{itemize}
    \item \textbf{No multicolinealidad:} Las variables predictoras no están perfectamente correlacionadas entre sí
\end{itemize}

\subsection{Interpretación de coeficientes}

En regresión múltiple, cada $\beta_j$ se interpreta como el cambio esperado en $Y$ cuando $X_j$ aumenta en 1 unidad, \textbf{manteniendo constantes} las demás variables.

\textbf{Ejemplo:}
\[
\hat{Y} = 50 + 1.2 X_1 + 0.8 X_2
\]
donde $Y$ = precio de casa (miles \$), $X_1$ = m$^2$, $X_2$ = número de habitaciones.

\begin{itemize}
    \item $\hat{\beta}_1 = 1.2$: Por cada m$^2$ adicional, el precio aumenta 1.2 mil \$, \textbf{para un número fijo de habitaciones}
    \item $\hat{\beta}_2 = 0.8$: Por cada habitación adicional, el precio aumenta 0.8 mil \$, \textbf{para un tamaño fijo}
\end{itemize}

\subsection{$R^2$ ajustado}

El $R^2$ siempre aumenta al agregar variables, aunque no aporten información. El \textbf{$R^2$ ajustado} penaliza la complejidad:

\begin{equation}
R^2_{\text{adj}} = 1 - \frac{\text{SCR}/(n - p - 1)}{\text{SCT}/(n - 1)} = 1 - (1 - R^2)\frac{n - 1}{n - p - 1}
\end{equation}

$R^2_{\text{adj}}$ puede disminuir si se agrega una variable irrelevante.

\subsection{Test F global}

Para probar si \textbf{al menos una} variable es útil:
\[
H_0: \beta_1 = \beta_2 = \cdots = \beta_p = 0 \quad \text{vs} \quad H_1: \text{al menos un } \beta_j \neq 0
\]

Estadístico:
\begin{equation}
F = \frac{\text{SCE}/p}{\text{SCR}/(n - p - 1)} \sim F_{p, n-p-1}
\end{equation}

Si $F > F_{p, n-p-1, \alpha}$, rechazamos $H_0$.

\subsection{Tests t individuales}

Para cada coeficiente $\beta_j$:
\[
H_0: \beta_j = 0 \quad \text{vs} \quad H_1: \beta_j \neq 0
\]

\begin{equation}
t_j = \frac{\hat{\beta}_j}{\text{SE}(\hat{\beta}_j)} \sim t_{n-p-1}
\end{equation}

\textbf{Importante:} Un test t significativo indica que $X_j$ aporta información \textbf{dado que las demás variables ya están en el modelo}.

\subsection{Multicolinealidad}

\textbf{Definición:} Ocurre cuando dos o más predictores están altamente correlacionados.

\textbf{Consecuencias:}
\begin{itemize}
    \item Estimaciones de $\beta_j$ inestables (alta varianza)
    \item Interpretación de coeficientes difícil
    \item Tests t individuales no significativos, pero test F sí
\end{itemize}

\textbf{Detección:}
\begin{itemize}
    \item \textbf{VIF (Factor de Inflación de la Varianza):}
    \[
    \text{VIF}_j = \frac{1}{1 - R^2_j}
    \]
    donde $R^2_j$ es el $R^2$ de la regresión de $X_j$ sobre las demás variables.
    
    \textbf{Regla práctica:} $\text{VIF} > 10$ indica multicolinealidad problemática
\end{itemize}

\textbf{Soluciones:}
\begin{itemize}
    \item Eliminar variables redundantes
    \item Combinar variables correlacionadas
    \item Usar regresión ridge o lasso (regularización)
    \item Análisis de componentes principales
\end{itemize}

%=============================================================================
\section{Selección de Variables}
%=============================================================================

\subsection{¿Por qué seleccionar variables?}

\textbf{Ventajas de modelos más simples:}
\begin{itemize}
    \item Más interpretables
    \item Menor riesgo de sobreajuste
    \item Predicciones más estables
    \item Menor costo computacional
\end{itemize}

\subsection{Criterios de selección}

\textbf{1. AIC (Criterio de Información de Akaike):}
\begin{equation}
\text{AIC} = n \ln(\text{SCR}/n) + 2(p + 1)
\end{equation}

\textbf{2. BIC (Criterio de Información Bayesiano):}
\begin{equation}
\text{BIC} = n \ln(\text{SCR}/n) + \ln(n)(p + 1)
\end{equation}

BIC penaliza más la complejidad que AIC.

\textbf{Regla:} Menor AIC o BIC $\Rightarrow$ mejor modelo

\subsection{Métodos de selección}

\textbf{1. Selección hacia adelante (Forward):}
\begin{itemize}
    \item Comenzar con el modelo vacío (solo intercepto)
    \item Agregar la variable que más mejora el criterio
    \item Repetir hasta que ninguna mejora el criterio
\end{itemize}

\textbf{2. Eliminación hacia atrás (Backward):}
\begin{itemize}
    \item Comenzar con el modelo completo (todas las variables)
    \item Eliminar la variable menos significativa
    \item Repetir hasta que todas sean significativas
\end{itemize}

\textbf{3. Selección por pasos (Stepwise):}
\begin{itemize}
    \item Combina forward y backward
    \item En cada paso, puede agregar o eliminar variables
\end{itemize}

\textbf{Advertencia:} Estos métodos son heurísticos y pueden no encontrar el mejor modelo. Es recomendable combinarlos con conocimiento del dominio.

%=============================================================================
\section{Regresión Polinomial}
%=============================================================================

\subsection{Motivación}

Cuando la relación entre $X$ e $Y$ no es lineal, podemos ajustar un polinomio:

\begin{equation}
Y_i = \beta_0 + \beta_1 X_i + \beta_2 X_i^2 + \cdots + \beta_d X_i^d + \epsilon_i
\end{equation}

donde $d$ es el grado del polinomio.

\textbf{Nota:} Aunque la relación con $X$ no es lineal, el modelo sigue siendo \textbf{lineal en los parámetros} $\beta_j$, por lo que se puede usar MCO.

\subsection{Implementación}

Crear nuevas variables:
\begin{itemize}
    \item $X_1 = X$
    \item $X_2 = X^2$
    \item $X_3 = X^3$, etc.
\end{itemize}

Luego ajustar una regresión lineal múltiple estándar.

\subsection{Selección del grado $d$}

\textbf{Métodos:}
\begin{itemize}
    \item Validación cruzada
    \item Criterios de información (AIC, BIC)
    \item Análisis visual de ajuste
\end{itemize}

\textbf{Advertencia:} Polinomios de grado muy alto pueden sobreajustar los datos.

%=============================================================================
\section{Correlación}
%=============================================================================

\subsection{Correlación de Pearson}

Mide la \textbf{fuerza y dirección} de la relación lineal entre dos variables $X$ e $Y$:

\begin{equation}
r_{XY} = \frac{\sum_{i=1}^{n}(X_i - \bar{X})(Y_i - \bar{Y})}{\sqrt{\sum_{i=1}^{n}(X_i - \bar{X})^2} \sqrt{\sum_{i=1}^{n}(Y_i - \bar{Y})^2}} = \frac{S_{XY}}{\sqrt{S_{XX}S_{YY}}}
\end{equation}

\textbf{Propiedades:}
\begin{itemize}
    \item $-1 \leq r_{XY} \leq 1$
    \item $r_{XY} = 1$: correlación positiva perfecta
    \item $r_{XY} = -1$: correlación negativa perfecta
    \item $r_{XY} = 0$: no hay correlación lineal (pero puede haber relación no lineal)
\end{itemize}

\textbf{Interpretación:}
\begin{itemize}
    \item $|r| < 0.3$: correlación débil
    \item $0.3 \leq |r| < 0.7$: correlación moderada
    \item $|r| \geq 0.7$: correlación fuerte
\end{itemize}

\subsection{Test de significancia}

Para probar $H_0: \rho = 0$ vs $H_1: \rho \neq 0$:

\begin{equation}
t = \frac{r \sqrt{n - 2}}{\sqrt{1 - r^2}} \sim t_{n-2}
\end{equation}

\textbf{Supuestos:}
\begin{itemize}
    \item $(X, Y)$ siguen una distribución normal bivariada
    \item Relación lineal
\end{itemize}

\subsection{Correlación de Spearman}

\textbf{Uso:} Cuando los datos no son normales o la relación es monótona pero no lineal.

\textbf{Método:} Calcular Pearson sobre los \textbf{rangos} de $X$ e $Y$:

\begin{equation}
r_S = 1 - \frac{6 \sum_{i=1}^{n} d_i^2}{n(n^2 - 1)}
\end{equation}

donde $d_i$ es la diferencia entre los rangos de $X_i$ e $Y_i$.

\textbf{Ventajas:}
\begin{itemize}
    \item Robusto a outliers
    \item No requiere normalidad
    \item Detecta relaciones monótonas no lineales
\end{itemize}

\subsection{Correlación no implica causalidad}

\textbf{Importante:} Una correlación significativa entre $X$ e $Y$ NO implica que $X$ cause $Y$ (o viceversa).

\textbf{Posibles explicaciones:}
\begin{itemize}
    \item $X$ causa $Y$
    \item $Y$ causa $X$
    \item Ambas están influenciadas por una tercera variable $Z$ (variable confusora)
    \item Correlación espuria (coincidencia)
\end{itemize}

\textbf{Ejemplo clásico:} Alta correlación entre ventas de helados y ahogamientos. No es que los helados causen ahogamientos, sino que ambos aumentan en verano.

%=============================================================================
\section{Análisis de Covarianza (ANCOVA)}
%=============================================================================

\subsection{Definición y objetivo}

ANCOVA combina ANOVA con regresión lineal. Permite comparar grupos (factor categórico) mientras se ajusta por variables continuas (covariables).

\textbf{Modelo general:}
\begin{equation}
Y_{ij} = \mu + \alpha_i + \beta X_{ij} + \epsilon_{ij}
\end{equation}

donde:
\begin{itemize}
    \item $Y_{ij}$: variable respuesta para la $j$-ésima observación en el grupo $i$
    \item $\alpha_i$: efecto del grupo $i$
    \item $X_{ij}$: covariable continua
    \item $\beta$: pendiente común (supuesto de pendientes paralelas)
    \item $\epsilon_{ij} \sim N(0, \sigma^2)$
\end{itemize}

\subsection{¿Cuándo usar ANCOVA?}

\textbf{Situaciones típicas:}
\begin{itemize}
    \item Comparar tratamientos ajustando por una medida basal (pre-test)
    \item Reducir error experimental controlando por variables conocidas
    \item Aumentar la potencia estadística
\end{itemize}

\textbf{Ejemplo:} Comparar tres dietas ($A$, $B$, $C$) sobre la pérdida de peso, ajustando por peso inicial.

\subsection{Supuestos}

Además de los supuestos de regresión lineal:

\begin{enumerate}
    \item \textbf{Independencia de la covariable y el tratamiento:} $X$ no debe estar influenciada por el grupo
    \item \textbf{Pendientes paralelas (homogeneidad de regresión):} La relación entre $X$ e $Y$ es la misma en todos los grupos
    \item \textbf{Linealidad:} La relación entre la covariable y la respuesta es lineal
\end{enumerate}

\subsection{Test de pendientes paralelas}

Ajustar el modelo con interacción:
\begin{equation}
Y_{ij} = \mu + \alpha_i + \beta X_{ij} + (\alpha \beta)_i X_{ij} + \epsilon_{ij}
\end{equation}

\textbf{Hipótesis:}
\[
H_0: (\alpha\beta)_1 = (\alpha\beta)_2 = \cdots = 0 \quad \text{(pendientes iguales)}
\]

Si no rechazamos $H_0$, procedemos con ANCOVA. Si la rechazamos, las pendientes difieren y ANCOVA no es apropiado.

\subsection{Interpretación de resultados}

Después de ajustar por la covariable:
\begin{itemize}
    \item \textbf{Efecto de la covariable:} $\hat{\beta}$ cuantifica cuánto cambia $Y$ por unidad de $X$
    \item \textbf{Efecto del grupo:} $\hat{\alpha}_i$ son las diferencias entre grupos \textbf{ajustadas} por $X$
\end{itemize}

\textbf{Ventaja principal:} ANCOVA reduce el error residual al explicar parte de la variabilidad con la covariable, aumentando la potencia para detectar diferencias entre grupos.

%=============================================================================
\section{Implementación en Python y R}
%=============================================================================

\subsection{Regresión lineal simple en Python}

\begin{lstlisting}[language=Python, caption=Regresión lineal simple con statsmodels]
import pandas as pd
import numpy as np
import matplotlib.pyplot as plt
import statsmodels.api as sm
from statsmodels.formula.api import ols

# Datos de ejemplo
data = pd.DataFrame({
    'X': [1, 2, 3, 4, 5, 6, 7, 8, 9, 10],
    'Y': [2.1, 4.3, 5.8, 8.2, 10.1, 11.8, 14.2, 15.9, 18.1, 20.3]
})

# Ajustar modelo
model = ols('Y ~ X', data=data).fit()

# Resumen del modelo
print(model.summary())

# Coeficientes
print(f"Intercepto: {model.params['Intercept']:.3f}")
print(f"Pendiente: {model.params['X']:.3f}")
print(f"R^2: {model.rsquared:.3f}")

# Grafico
plt.figure(figsize=(8, 5))
plt.scatter(data['X'], data['Y'], label='Datos')
plt.plot(data['X'], model.fittedvalues, 'r-', label='Ajuste')
plt.xlabel('X')
plt.ylabel('Y')
plt.legend()
plt.show()
\end{lstlisting}

\subsection{Diagnóstico en Python}

\begin{lstlisting}[language=Python, caption=Gráficos de diagnóstico]
import matplotlib.pyplot as plt
import scipy.stats as stats

fig, axes = plt.subplots(2, 2, figsize=(12, 10))

# 1. Residuos vs Valores Ajustados
axes[0, 0].scatter(model.fittedvalues, model.resid)
axes[0, 0].axhline(0, color='red', linestyle='--')
axes[0, 0].set_xlabel('Valores Ajustados')
axes[0, 0].set_ylabel('Residuos')
axes[0, 0].set_title('Residuos vs Ajustados')

# 2. Q-Q Plot
stats.probplot(model.resid, dist="norm", plot=axes[0, 1])
axes[0, 1].set_title('Q-Q Plot')

# 3. Scale-Location
residuos_std = np.sqrt(np.abs(model.resid_pearson))
axes[1, 0].scatter(model.fittedvalues, residuos_std)
axes[1, 0].set_xlabel('Valores Ajustados')
axes[1, 0].set_ylabel('Raiz(|Residuos Estand.|)')
axes[1, 0].set_title('Scale-Location')

# 4. Residuos vs Leverage
influence = model.get_influence()
leverage = influence.hat_matrix_diag
axes[1, 1].scatter(leverage, model.resid_pearson)
axes[1, 1].set_xlabel('Leverage')
axes[1, 1].set_ylabel('Residuos Estandarizados')
axes[1, 1].set_title('Residuos vs Leverage')

plt.tight_layout()
plt.show()
\end{lstlisting}

\subsection{Regresión múltiple en R}

\begin{lstlisting}[language=R, caption=Regresión múltiple en R]
# Cargar datos
data(mtcars)

# Ajustar modelo
modelo <- lm(mpg ~ wt + hp + cyl, data = mtcars)

# Resumen
summary(modelo)

# Coeficientes
coef(modelo)

# R^2 y R^2 ajustado
summary(modelo)$r.squared
summary(modelo)$adj.r.squared

# ANOVA
anova(modelo)

# Intervalos de confianza
confint(modelo)

# Prediccion
nuevos_datos <- data.frame(wt = 3.0, hp = 150, cyl = 6)
predict(modelo, nuevos_datos, interval = "prediction")
\end{lstlisting}

\subsection{Selección de variables en R}

\begin{lstlisting}[language=R, caption=Selección stepwise]
# Modelo completo
modelo_completo <- lm(mpg ~ ., data = mtcars)

# Seleccion stepwise (ambas direcciones)
modelo_step <- step(modelo_completo, direction = "both")

# Resumen del modelo seleccionado
summary(modelo_step)

# AIC y BIC
AIC(modelo_step)
BIC(modelo_step)
\end{lstlisting}

\subsection{ANCOVA en Python}

\begin{lstlisting}[language=Python, caption=ANCOVA con statsmodels]
import pandas as pd
from statsmodels.formula.api import ols
from statsmodels.stats.anova import anova_lm

# Datos de ejemplo: efecto de dieta sobre peso, ajustando por peso inicial
data = pd.DataFrame({
    'peso_final': [72, 68, 75, 70, 65, 78, 71, 69, 73, 67],
    'peso_inicial': [80, 78, 85, 82, 75, 88, 81, 79, 84, 76],
    'dieta': ['A', 'A', 'A', 'B', 'B', 'B', 'C', 'C', 'C', 'C']
})

# Modelo ANCOVA
modelo_ancova = ols('peso_final ~ C(dieta) + peso_inicial', 
                    data=data).fit()

# Tabla ANOVA
print(anova_lm(modelo_ancova, typ=2))

# Resumen
print(modelo_ancova.summary())

# Test de pendientes paralelas (interaccion)
modelo_interaccion = ols('peso_final ~ C(dieta) * peso_inicial', 
                         data=data).fit()
print(anova_lm(modelo_interaccion, typ=2))
\end{lstlisting}

%=============================================================================
\section{Ejercicios Propuestos}
%=============================================================================

\subsection{Ejercicio 1: Regresión lineal simple}

Se midió la altura ($X$, en cm) y el peso ($Y$, en kg) de 15 estudiantes. Los datos se resumen:
\[
\bar{X} = 170, \quad \bar{Y} = 68, \quad S_{XX} = 450, \quad S_{YY} = 280, \quad S_{XY} = 330
\]

\textbf{Tareas:}
\begin{enumerate}
    \item Calcular las estimaciones $\hat{\beta}_0$ y $\hat{\beta}_1$
    \item Interpretar los coeficientes
    \item Calcular $R^2$
    \item Predecir el peso de un estudiante de 175 cm
\end{enumerate}

\subsection{Ejercicio 2: Diagnóstico}

Ajustaste una regresión lineal y observaste en el gráfico de residuos vs ajustados un patrón en forma de embudo.

\textbf{Preguntas:}
\begin{enumerate}
    \item ¿Qué supuesto se viola?
    \item ¿Qué consecuencias tiene?
    \item Proponer dos posibles soluciones
\end{enumerate}

\subsection{Ejercicio 3: Regresión múltiple}

Se ajustó un modelo de regresión múltiple para predecir el precio de autos usados ($Y$) en función de:
\begin{itemize}
    \item $X_1$: antigüedad (años)
    \item $X_2$: kilometraje (miles)
    \item $X_3$: potencia del motor (HP)
\end{itemize}

El modelo ajustado es:
\[
\hat{Y} = 25000 - 1500 X_1 - 0.05 X_2 + 100 X_3
\]

\textbf{Tareas:}
\begin{enumerate}
    \item Interpretar cada coeficiente
    \item Si un auto tiene 3 años, 50,000 km y 150 HP, ¿cuál es su precio estimado?
    \item ¿Qué variable tiene mayor impacto absoluto en el precio?
\end{enumerate}

\subsection{Ejercicio 4: Multicolinealidad}

Se ajustó un modelo con las variables: estatura ($X_1$), peso ($X_2$) y circunferencia de cintura ($X_3$) para predecir presión arterial. Los VIF son:
\[
\text{VIF}_1 = 2.3, \quad \text{VIF}_2 = 15.2, \quad \text{VIF}_3 = 14.8
\]

\textbf{Preguntas:}
\begin{enumerate}
    \item ¿Hay multicolinealidad? ¿Entre qué variables?
    \item ¿Qué harías para solucionarlo?
\end{enumerate}

\subsection{Ejercicio 5: Correlación}

Se calculó la correlación de Pearson entre horas de estudio ($X$) y calificación en un examen ($Y$) para 50 estudiantes, obteniendo $r = 0.72$.

\textbf{Tareas:}
\begin{enumerate}
    \item Interpretar el valor de $r$
    \item Realizar el test de significancia ($\alpha = 0.05$)
    \item Calcular $R^2$ y explicar su significado
    \item ¿Podemos concluir que estudiar más causa mejores notas? Justificar
\end{enumerate}

%=============================================================================
\section{Referencias}
%=============================================================================

\begin{itemize}
    \item Kutner, M. H., Nachtsheim, C. J., Neter, J., \& Li, W. (2005). \textit{Applied Linear Statistical Models} (5th ed.). McGraw-Hill/Irwin.
    \item James, G., Witten, D., Hastie, T., \& Tibshirani, R. (2021). \textit{An Introduction to Statistical Learning with Applications in R} (2nd ed.). Springer.
    \item Montgomery, D. C., Peck, E. A., \& Vining, G. G. (2012). \textit{Introduction to Linear Regression Analysis} (5th ed.). Wiley.
    \item Fox, J. (2015). \textit{Applied Regression Analysis and Generalized Linear Models} (3rd ed.). SAGE Publications.
    \item Faraway, J. J. (2014). \textit{Linear Models with R} (2nd ed.). CRC Press.
\end{itemize}

\end{document}
